\documentclass[11pt]{article}
\usepackage[T1]{fontenc}
\usepackage[utf8]{inputenc}
\usepackage{lmodern,amsmath,amssymb,amsthm,mathtools,booktabs,array}
\usepackage[a4paper,margin=27mm]{geometry}
\usepackage[colorlinks=true,allcolors=blue]{hyperref}
\usepackage{microtype}
\hypersetup{pdftitle={Sharp growth of vertex-position ratios and realizable pairs},
pdfauthor={Haobo Ma; Wenlin Zhang}}
\newtheorem{theorem}{Theorem}
\newtheorem{lemma}[theorem]{Lemma}
\newtheorem{proposition}[theorem]{Proposition}
\newtheorem{corollary}[theorem]{Corollary}
\theoremstyle{remark}\newtheorem{remark}[theorem]{Remark}
\newcommand{\vp}{\operatorname{vp}}
\newcommand{\vpm}{\operatorname{vp}^{-}}
\newcommand{\dist}{\operatorname{dist}}
\newcommand{\diam}{\operatorname{diam}}
\newcommand{\vv}{\operatorname{vv}}
\newcommand{\vvm}{\operatorname{vv}^{-}}
\newcommand{\Z}{\mathbb Z}
\DeclareMathOperator*{\argmin}{arg\,min}
\title{Sharp growth of vertex-position ratios\newline and realizable pairs}
\author{Haobo Ma\\\small ChronoAI Pte Ltd; The Omega Institute
\and Wenlin Zhang\\\small National University of Singapore; The Omega Institute}
\date{13 October 2026}
\begin{document}
\maketitle
\begin{abstract}
We prove that the largest ratio of the upper to lower vertex-position number
among connected cubic graphs of order $n$ is $\Theta(n/\log n)$, for all
sufficiently large even $n$. For connected graphs of maximum degree exactly
$\Delta$, the corresponding maximum is $\Theta(n/(\Delta+\log n))$, uniformly
for $3\le\Delta\le n-1$. The proof extends a shortest-path-cover argument to
geodesic antichains and constructs cubic graphs with linearly many incomparable
vertices and logarithmically many root geodesics covering the graph. We also
classify all pairs $(\vpm,\vp)$ with lower parameter at most two: they are
$(1,1),(1,2)$ and $(2,b)$ for $2\le b\le10$; the largest upper parameter in
row $a$ grows exponentially in $a$. Finally, vertex-visibility ratios grow
linearly along a cubic family, and local monophonic-position ratios are
unbounded.
\end{abstract}

\section{Vertex position and the main results}
All graphs are finite, simple, undirected and connected. Following
Thankachy et al.\ \cite{position}, an $x$-position set is a set $S\subseteq V(G)$
such that, for each $v\in S$, no vertex of $S\setminus\{v\}$ lies on any
shortest $x,v$-path. Write $p_x(G)$ for its maximum size, and put
\[
 a=\vpm(G)=\min_xp_x(G),\qquad b=\vp(G)=\max_xp_x(G),\qquad \rho(G)=b/a.
\]
For $|V(G)|\ge2$ a maximum set can be chosen to exclude $x$; a set containing
$x$ has size at most one. We use $p_x(K_1)=1$.

Problem 2.4 of \cite{position} asks whether $\rho$ is bounded over connected
graphs. James Tuite asked in addition about bounded degree, local path
variants, and realizable pairs (private communication, 11 October 2026).
The following results give the sharp order of growth and settle the first
two rows of the pair problem. All logarithms in estimates below have base two.

\begin{theorem}\label{thm:main}
There are absolute positive constants bounding the following maxima above
and below by constant multiples of the displayed quantities:
\[
\begin{array}{ll}
\displaystyle\max_{|V(G)|=n,\ G\text{ cubic}}\rho(G)
   &\displaystyle=\Theta\!\left(\frac n{\log n}\right),\quad n\text{ even and sufficiently large},\\[6pt]
\displaystyle\max_{|V(G)|=n,\ \Delta(G)\le\Delta}\rho(G)
   &\displaystyle=\Theta\!\left(\frac n{\log n}\right),\quad \Delta\ge3,\ n\text{ sufficiently large},\\[6pt]
\displaystyle\max_{|V(G)|=n,\ \Delta(G)=\Delta}\rho(G)
   &\displaystyle=\Theta\!\left(\frac n{\Delta+\log n}\right),\quad 3\le\Delta\le n-1.
\end{array}
\]
In the last line the constants and the threshold for $n$ are independent of
$\Delta$. If $a\le2$, the realizable pairs $(a,b)$ are exactly
\[
 (1,1),\quad(1,2),\quad(2,2),(2,3),\ldots,(2,10).
\]
For every $a\ge2$, the maximum $B(a)=\max\{\vp(G):\vpm(G)=a\}$ exists and satisfies
\begin{equation}\label{eq:Ba}
 2^{a-1}\le B(a)\le(2a+1)4^a.
\end{equation}
\end{theorem}

For maximum degree at most two, paths and cycles give a maximum ratio of two
once $n\ge3$. Trees have ratio at most two: if a tree has $\ell\ge2$ leaves,
then $p_x=\ell-1$ at a leaf and $p_x=\ell$ elsewhere, apart from $K_2$
\cite[Corollary 3.6]{position}. The constructions below therefore require
cycles even though their low-degree branching comes from trees.

\section{An upper bound from geodesic covers}
Fix $x$ and define the geodesic order on $V(G)\setminus\{x\}$ by
\[
 u\preceq_x v\quad\Longleftrightarrow\quad d(x,u)+d(u,v)=d(x,v).
\]
The triangle inequality proves transitivity, and strict order increases
distance from $x$, so this is a partial order. Its antichains are exactly
the observer-excluding $x$-position sets. A sphere about $x$ is one such set.

\begin{lemma}\label{lem:rootcover}
If $p_y(G)=k$, then $V(G)$ is covered by $k$ geodesics starting at $y$.
Conversely, a cover by $k$ such geodesics implies $p_y(G)\le k$.
\end{lemma}
\begin{proof}
Dilworth's theorem \cite{dilworth} partitions the nonroot vertices into $k$
chains. For a chain $v_1\prec_y\cdots\prec_y v_s$, concatenate a shortest
$y,v_1$-path and shortest paths between successive chain vertices. The total
length telescopes to $d(y,v_s)$, so the concatenation is a geodesic containing
the chain. Conversely, a position set meets each rooted geodesic in at most
one vertex.
\end{proof}

We need a bound on all antichains, including those occupying several spheres.
The following proof gives the endpointwise extension of the vertex-cover
construction of Dumas--Foucaud--Perez--Todinca
\cite[Section 4, proof of Lemma 8]{cover}. Their recursion bounds spheres;
the final slack argument below replaces the equal-radius condition.

\begin{lemma}\label{lem:coverbound}
If $V(G)$ is covered by $k$ isometric paths, then
\[
 p_x(G)\le(2k+1)4^k\qquad\text{for every }x.
\]
\end{lemma}
\begin{proof}
Orient the covering paths $\mu_1,\ldots,\mu_k$ and call their indices colors.
A vertex may receive the color of any base path containing it. Call a colored
walk good if each color occurs in one block, following its base path
monotonically. Every pair of vertices has a good shortest path: among colored
shortest paths minimize the number of blocks; if a color occurs twice,
replace the intervening part by the isometric base segment and color it
uniformly, contradicting minimality. A block has sign $+$ or $-$ according
to its direction; assign $+$ to a singleton block.

Take good shortest $x,v$-paths for $v$ in an $x$-position set $S$. We describe
the recursive grouping used in \cite{cover}, recording lengths separately
for each endpoint. At each node there is one current good walk to each
endpoint in its group. Stop if all these walks are colored prefixes of one
of them. Initially let $C$ be the colors assigned to $x$. In each walk take
the last vertex $w$ colored in $C$, of color $c$, and replace the prefix to
$w$ by $\mu_c[x,w]$. Group the resulting walks by $(c,\text{sign})$.
This first replacement does not increase length.

At a subsequent nonterminal node, take a longest colored prefix $R$ such that
every walk either extends $R$ or is a prefix of $R$. Let $z$ be its last vertex
and set aside the walks ending within $R$. Call the others unfinished.
Let $C$ be the colors immediately after $z$. In each unfinished
walk take the last vertex $w$ whose color $c$ lies in $C$. Replace its part
after $z$ through $w$ by one edge from $z$ to a vertex $u$ of $\mu_c$, followed
by $\mu_c[u,w]$. If $z\in\mu_c$, take $u$ toward $w$ along that path.
Otherwise choose the first neighbor of $z$ in the fixed orientation of
$\mu_c$. Such a neighbor exists because $c\in C$.
Group by the color and direction of this replacement. If the new block is
a singleton and the color of $z$ is also $c$, use the direction of $zu$.
Assign the previously ended prefixes to any one nonempty group.

Here are the invariants that control this recursion. The last-occurrence
choice removes all other colors in $C$ from the continuing suffix. If $c$
occurred in the retained prefix, it was its last color, so the replacement
still has one monotone block of each color. Walks in the same group have the
same successor of $z$. Thus their retained prefix advances at the next
nonterminal node. After the first grouping, the active suffixes have a common
initial color and direction. At a split with $|C|\ge2$, the chosen child loses
all colors of $C$ except its chosen color. At a split with $C=\{c\}$, the
previous initial color is different from $c$: otherwise all continuing walks
would share the next base-path edge, contrary to maximality of the common
prefix. That previous color has been passed and cannot recur. Consequently
the union of available suffix colors strictly decreases at every subsequent
split. There are at most $k$ grouping steps from the root to a terminal node.

Children have distinct labels in the alphabet of $2k$ color--sign pairs.
In any child, all the other sibling labels are forbidden in its descendants:
the other colors have been removed; the opposite sign of its own color
cannot recur in a good walk. These statements also hold for ended prefixes,
which are carried unchanged and never create a branch.

Every replacement adds at most two edges. This is immediate if $z\in\mu_c$.
Otherwise isometry gives
\[
 |\mu_c[u,w]|=d(u,w)\le1+|W[z,w]|,
 \qquad 1+|\mu_c[u,w]|\le |W[z,w]|+2,
\]
where $W$ is the walk being modified. Thus at recursion depth $j$ the walk
to endpoint $v$ has length at most $d(x,v)+2j$. This estimate uses no common
endpoint radius. It is also valid if replacement walks repeat vertices:
the coloring refers to occurrences and the retained-prefix argument applies
to walks.

The number of terminal groups is at most $4^k$. More generally, with $q$
available labels there are at most $2^q$ terminal descendants. Induct on the
remaining tree height. At a node with $t\ge2$ children each child has at most
$q-t+1$ available labels, and
\[
 t\,2^{q-t+1}\le2^q.
\]
A unary node preserves the bound. Initially $q=2k$.

In one terminal group all endpoint walks are prefixes of a common walk $W$.
Write $\ell_v$ for the prefix length and $e_v=\ell_v-d(x,v)$ for its slack.
Then $0\le e_v\le2k$. If $u$ precedes $v$ on $W$, the triangle inequality
gives $e_u\le e_v$. If equality holds, then
\[
 d(x,v)-d(x,u)=\ell_v-\ell_u\ge d(u,v)
                         \ge d(x,v)-d(x,u),
\]
so $u\preceq_x v$, contrary to the antichain condition. The integer slacks
therefore strictly increase along the endpoints. Each terminal group has
at most $2k+1$ endpoints, proving the result.
\end{proof}

\begin{corollary}\label{cor:upper}
Every connected graph of order $n\ge2$ satisfies
\begin{equation}\label{eq:universal}
 b\le(2a+1)4^a,\qquad \rho(G)\le\frac{4n}{\log n}.
\end{equation}
\end{corollary}
\begin{proof}
Choose $y$ with $p_y=a$ and apply Lemmas~\ref{lem:rootcover} and
\ref{lem:coverbound}. Hence
\[
 \rho\le\min\{(n-1)/a,\,3\cdot4^a\}.
\]
For $a\ge\frac14\log n$ use the first term. Otherwise the second is less than
$3\sqrt n\le4n/\log n$. The last inequality follows, for example, by maximizing
$(\log_2 t)/\sqrt t$ over positive real $t$, whose maximum is $2/(e\ln2)<4/3$.
\end{proof}

\section{Cubic graphs attaining the growth order}
We first build a subcubic graph with many incomparable targets. Its dyadic
part is related to the exponential-sphere construction in
\cite[Figure 7]{cover}; a periodic binary-tree arrangement replicates its
targets across enough rays to make their number proportional to the order.

\subsection{The base graph}
Fix $h\ge8$, let $q=2^h$ and $d=\lfloor\log h\rfloor$, and take a full binary
tree $T$ of depth $d$, rooted at $r$. Write
\[
 \ell=2^d,\quad R=2^{d+1}-1,\quad w=h+R,
\]
for its number of leaves, number of vertices, and the required number of rays.
Choose a proper edge coloring $\chi:E(T)\to\{1,2,3\}$; recursively assign the
two child edges colors different from the parent-edge color. Put
\begin{align*}
 A&=4(q-1),& M&=3d+6,& F&=A+h,& K&=F+M+1,\\
 t_0&=8\left\lceil\frac{h+M+2}{8}\right\rceil,
 &s&=A+t_0,&L&=t_0+8q+8M.
\end{align*}
Let $k$ be the least power of two at least $w$. Take a full binary hub tree
with $k$ leaves and root $y$. Attach a length-$L$ ray to each of $w$ distinct
hub leaves, indexed from level $1$ beyond its hub leaf. Name the rays
$P_0,\ldots,P_{h-1}$ and $U_u$ for $u\in V(T)$, with $P_h=U_r$.
Define
\[
 \delta_i=4\cdot2^{h-i},\qquad C_0=\{s\},\qquad
 C_i=\{c-\delta_i,c+\delta_i:c\in C_{i-1}\}.
\]
The last set is
\begin{equation}\label{eq:centers}
 C_h=\{t_0+8j:0\le j<q\}.
\end{equation}
Add precisely the edges
\begin{align}
 P_{i-1,t}P_{i,t}&\quad(1\le i\le h,\ t\in C_i),\label{eq:dyadicedges}\\
 U_{u,t}U_{v,t}&\quad(uv\in E(T),\ 1\le t\le L,\ t-t_0\equiv\chi(uv)\pmod8).\label{eq:treeedges}
\end{align}
Call this graph $B^{(h)}$ and choose $x=P_{0,s}$.

The sets $C_i$ and $C_{i+1}$ are disjoint: their elements have different
residues modulo $2\delta_{i+1}$. Thus an intermediate generator ray receives
at most one extra edge at each vertex. At $U_r$ the incoming dyadic edges
have residue $0$ modulo $8$, whereas the tree edges have residues $1,2,3$.
Proper edge coloring gives the same separation at every other $U$-ray.
The added edges avoid the ray endpoints, since $t_0\ge2$ and $L\equiv0\pmod8$
lies beyond all dyadic centers. Hence the maximum degree is three, the root
has degree two, and there are exactly $k$ leaves, including unused hub leaves.

Every added edge preserves ray level. The hub-root paths, extended along
their attached rays, remain geodesics and cover the graph. The $k$ leaves
form a $y$-position set, so
\begin{equation}\label{eq:baseorder}
 p_y(B^{(h)})=k,\qquad N_h=|V(B^{(h)})|=2k-1+wL=\Theta(h2^h).
\end{equation}
Here $h/2\le\ell\le h$, $w=\Theta(h)$, $k=\Theta(h)$ and $L=\Theta(2^h)$.

\subsection{Distances and the target antichain}
Let $H$ be the subgraph formed by the $U$-rays and
\eqref{eq:treeedges}, and let $C=\{U_{r,t}:t\in C_h\}$. Put
\[
 F_i=\sum_{j=1}^i(\delta_j+1),\qquad F_0=0,\qquad F_h=F.
\]
Define the potential
\begin{equation}\label{eq:potential}
 f(v)=
 \begin{cases}
 K,&v\text{ in the hub tree},\\
 \min\{K,F_i+\dist(t,C_i)\},&v=P_{i,t},\ 0\le i<h,\\
 \min\{K,F+d_H(C,v)\},&v\in V(H).
 \end{cases}
\end{equation}
It is $1$-Lipschitz on every edge. On rays and within $H$ this follows from
distance-to-a-set inequalities. On the root ray of $H$,
$d_H(C,U_{r,t})=\dist(t,C_h)$, since the level coordinate is $1$-Lipschitz
and the root ray realizes this lower bound. At a dyadic edge at level
$t\in C_i$, the nearest point of $C_{i-1}$ is at distance $\delta_i$;
the untruncated values on its two endpoints are $F_i-1$ and $F_i$.
To see the nearest-center assertion, centers in $C_{i-1}$ are separated
by at least $4\delta_i$, while each new center is $\delta_i$ from its parent.
Finally, the first vertex of generator ray $i$ has untruncated value
$s+i-1\ge K$, and the first vertex of a ray of $H$ has value at least
$F+t_0-1\ge K$. Thus edges to the hub also satisfy the bound.

Since $f(x)=0$, distances from $x$ are at least $f$. Inductively follow a
ray by $\delta_i$ steps from a parent center and cross one dyadic edge.
This gives length $F_i$ paths to all centers of $C_i$, and therefore
\begin{equation}\label{eq:centerdist}
 d(x,P_{i,t})=F_i\quad(t\in C_i),\qquad d(x,c)=F\quad(c\in C).
\end{equation}

Consider the finite quotient $Q=V(T)\times\Z_8$, with a cycle on each
$u$-fiber and the edge $(u,\chi(uv))(v,\chi(uv))$ for each $uv\in E(T)$.
For a leaf $u$ of $T$ let $c_u=d_Q((r,0),(u,0))$. Following its tree path,
using at most three cycle steps initially and finally and at most two
between successive edge colors, gives
\[
 c_u\le3+d+2(d-1)+3=3d+4<M.
\]
Set $m_0=\lceil M/8\rceil$ and $J=\{m_0,\ldots,q-1-m_0\}$. Then
\begin{equation}\label{eq:periodicdistance}
 d_H(C,U_{u,t_0+8j})=c_u\qquad(u\text{ a leaf},\ j\in J).
\end{equation}
Projection to $Q$ proves the lower bound. For the upper bound, lift a shortest
quotient path backwards from the indicated target. Its length is less than
$M$, so every lifted level differs from the target level by less than $M$.
The root-ray endpoint has residue $0$ and, by the margins defining $J$, lies
between $t_0$ and $t_0+8(q-1)$; hence it belongs to $C$. The same margins and
the padding in $L$ keep the entire lift inside $H$.

Combining \eqref{eq:potential}--\eqref{eq:periodicdistance} yields
\[
 d(x,U_{u,t_0+8j})=F+c_u<K.
\]
Let $S$ consist of all these targets. Targets on the same leaf ray have
equal $x$-distance. A geodesic to a target cannot meet the hub, where the
potential is $K$. Away from the hub, any visit to a different leaf ray must
enter and leave via its parent ray. If the entry and exit levels are $t,t'$,
the excursion has length at least $|t-t'|+2$ and can be shortened to the
parent-ray segment of length $|t-t'|$. Thus such a visit cannot occur in a
geodesic. It follows that $S$ is an $x$-position set.

For $h\ge8$, $m_0\le(3\log h+13)/8\le2^{h-2}$: the second inequality holds
at $h=8$ and its right side doubles at each step whereas the left side grows
by less than one. Consequently
\begin{equation}\label{eq:targetsize}
 |S|=\ell(q-2m_0)\ge\frac{hq}{4}.
\end{equation}

\subsection{Cubic completion and every sufficiently large order}
The following operation will preserve the target antichain.
\begin{lemma}\label{lem:cubic}
Let $B$ be a connected subcubic graph with a degree-two root $y$, exactly $k$
leaves, and a cover by $k$ geodesics from $y$ to those leaves. There is a
connected cubic graph $C(B)$ of order $2|V(B)|+4k$ containing an isometric
copy of $B$ and satisfying $p_{y^0}(C(B))\le3k$.
\end{lemma}
\begin{proof}
Take copies $B^0,B^1$. At every degree-two vertex add $v^0v^1$. At a leaf
$u$, add a fresh four-cycle $c_0c_1c_2c_3c_0$, joining $u^0$ to $c_0,c_1$
and $u^1$ to $c_2,c_3$. Every vertex now has degree three, and $y^0y^1$
connects the copies. Projection to $B$, collapsing each cycle to its leaf,
sends edges to edges or stationary steps. Thus each old copy is isometric.

Write $\lambda(v)=d_B(y,v)$. Distances from $y^0$ to $v^0,v^1$ are
$\lambda(v),\lambda(v)+1$. The first is isometry. For the second, a path
must switch copies, introducing at least one stationary step under projection,
while the root edge followed by a copy-one geodesic attains the bound.
At a leaf $u$ the four cycle distances are
$\lambda(u)+1,\lambda(u)+1,\lambda(u)+2,\lambda(u)+2$.
For each rooted leaf geodesic use three paths: its copy-zero version extended
by $c_0,c_3$; its copy-zero version extended by $c_1,c_2$; and $y^0y^1$
followed by its copy-one version. The distance formulas make all three
geodesics, and these $3k$ paths cover $C(B)$.
\end{proof}

Apply the lemma to $B^{(h)}$. Isometry preserves both observer distances
and pairwise distances, and hence preserves $S$ as a position set in copy zero.
The resulting graph $C^{(h)}$ has
\begin{equation}\label{eq:cubicorder}
 n_h=2N_h+4k=2wL+8k-2=\Theta(h2^h),\qquad
 \rho(C^{(h)})\ge\frac{\ell(q-2m_0)}{3k}=\Omega(2^h).
\end{equation}
The formulas for $N_h$ and $n_h$ increase with $h$; their two-sided
$\Theta(h2^h)$ estimates give bounded consecutive ratios. Extending a single ray beyond
all its attachments by $t$ vertices leaves old distances, the number of leaves,
and the size of the rooted geodesic cover unchanged. It adds $t$ vertices
to the base and $2t$ vertices to its cubic completion. Given any sufficiently
large $n$, choose the largest base order $N_h\le n$ and pad to $n$.
For an even $n$, choose the largest cubic order $n_h\le n$ and pad the base
by $(n-n_h)/2$. Bounded consecutive ratios keep the number of targets
proportional to $n$, while the cover size is $O(\log n)$.
This proves the lower bounds in the first two lines of Theorem~\ref{thm:main};
Corollary~\ref{cor:upper} proves their upper bounds.

For instance, $h=10$ gives $q=1024$, $d=3$, $\ell=8$, $w=25$, $k=32$,
$M=15$, $t_0=32$, and $L=8344$. Thus
\begin{equation}\label{eq:concrete}
 |V(C^{(10)})|=417454,\qquad \vp(C^{(10)})\ge8160,\qquad
 \vpm(C^{(10)})\le96,\qquad \rho(C^{(10)})\ge85.
\end{equation}

\subsection{Maximum degree exactly \texorpdfstring{$\Delta$}{Delta}}
Thankachy et al.\ \cite[Theorem 2.9]{position} prove
$a\ge\lceil(\Delta+1)/3\rceil$. Indeed, from any observer the closed
neighborhood of a maximum-degree vertex lies in three consecutive distance
layers; one is large enough to give the bound (if the observer is that
vertex, use its neighborhood). Together with \eqref{eq:universal}, this gives
\[
 \rho(G)\le\min\left\{\frac{3n}{\Delta+1},\frac{4n}{\log n}\right\}
        \le\frac{7n}{\Delta+1+\log n}.
\]
The final inequality follows by adding the two resulting lower bounds on
$1/\rho$ with weights $3/7$ and $4/7$.

For $3\le\Delta\le n/2$, take a padded base construction of order
$N=n-(\Delta-2)\ge n/2$ and attach $\Delta-2$ new leaves at its degree-two
root. Its maximum degree is exactly $\Delta$, old distances are preserved,
and its root cover has $k+\Delta-2=O(\log n+\Delta)$ paths. The old target
set has size $\Omega(N)=\Omega(n)$, with absolute constants. For
$\Delta>n/2$, the quantity $n/(\Delta+\log n)$ is bounded above and below
by absolute positive constants. A star with $\Delta$ leaves, with one edge
subdivided to obtain order $n$, has the required degree and ratio at least
one. This completes the degree statement.

\section{Realizable pairs}
If $p_y=1$, Lemma~\ref{lem:rootcover} gives a spanning geodesic. Any extra
edge would shortcut it, so the graph is a path. Conversely, a path endpoint
has value one. Thus the first row consists of $(1,1)$, attained by $K_1,K_2$,
and $(1,2)$, attained by every longer path, as in
\cite[Corollary 2.7]{position}.

\subsection{The exact second row}
\begin{proposition}\label{prop:ten}
If $\vpm(G)=2$, then $\vp(G)\le10$.
\end{proposition}
\begin{proof}
Choose $y$ with $p_y=2$ and cover the graph by two $y$-geodesics $P,Q$,
indexed by distance from $y$. They may share vertices. Put the observer at
$x=P_r$, interchanging the paths if necessary. Since $P$ is isometric,
$d(x,P_i)=|i-r|$, so an $x$-position set contains at most two vertices of $P$.

A contact is a shared vertex $P_i=Q_j$, of cost $\epsilon=0$, or an edge
$P_iQ_j$, of cost $\epsilon=1$. Shared vertices have $i=j$ and edges have
$|i-j|\le1$. Taking the last visit to $P$ in a shortest path gives
\begin{equation}\label{eq:contact}
 d(x,Q_t)=\min_{(i,j,\epsilon)}\{|i-r|+\epsilon+|t-j|\}.
\end{equation}
After its last visit the path lies in $Q$; no chord between nonconsecutive
vertices of $Q$ is possible because $Q$ is isometric. Conversely every
expression in \eqref{eq:contact} is the length of a walk. Write
\[
 |i-r|+\epsilon=|j-r|+c,\qquad c\in\{0,1,2\}.
\]
Separate contacts into $j\le r$ and $j\ge r$, assigning equality to one
group, and let $f_-,f_+$ be their lower envelopes. For a nonempty left group
let $J$ be its largest contact index. When $t\le J$,
\[
 f_-(t)=r-t+g(t),\qquad g(t)\in\{0,1,2\}.
\]
The lower bound is termwise and the contact at $J$ gives the upper bound.
Two selected vertices assigned to this envelope with the same $g$-value
are comparable along $Q$. There are at most three such selected vertices
with $t\le J$. For $t>J$ all terms have slope $+1$, contributing at most
one more. Reflection gives at most four for the right group as well.
Assign each selected vertex of $Q\setminus P$ to an envelope attaining
its actual distance; the comparisons therefore use actual distances.
There are at most eight selected vertices outside $P$, hence at most ten
in all.
\end{proof}

For explicit realizations we also need a bound that counts contacts.
\begin{lemma}\label{lem:rungs}
A graph made of two paths sharing just their root and $s$ additional rungs,
each joining positive levels differing by at most one, satisfies $p_x\le s+4$
for every observer $x$.
\end{lemma}
\begin{proof}
Use \eqref{eq:contact} with $x=P_r$. Suppose $q$ rungs are assigned to the
left group. If $q>0$, any one dominates the root contact on positive
$Q$-indices, because its offset is at most two and its level is at least one.
For these rungs
\[
 g(t)=f_-(t)+t-r=\min_j\{c_j+2\max(t-j,0)\}.
\]
Up to their last contact index $J$, $g(t)\in\{0,1,2\}$ and every attained
value is one of the offsets $c_j$: values zero and one require that offset;
if no offset is two, the last contact bounds $g\le1$ throughout this range.
There are thus at most $q$ selected vertices before or at $J$, and one in
the final slope-$+1$ tail. If $q=0$, the root contact contributes a single
chain. The reflected right group contributes at most its number of rungs
plus one (or zero if it is empty). The other ray therefore contributes at
most $s+2$, and the observer's ray at most two.
\end{proof}

Take two length-$13$ paths with common root $P_0=Q_0=y$ and define
\[
\begin{array}{lll}
 e_1=P_4Q_5,&e_2=P_3Q_3,&e_3=P_2Q_1,\\
 e_4=P_{10}Q_9,&e_5=P_{11}Q_{11},&e_6=P_{12}Q_{13}.
\end{array}
\]
The following table specifies five graphs and selected sets with observer
$x=P_7$. The last column lists the corresponding $x$-distances.
\begin{center}
\small
\begin{tabular}{@{}c l l l l@{}}
\toprule
$b$ & Rung indices & Selected on $P$ & Selected $Q$-indices & $Q$-distances\\
\midrule
6 & $1,4$ & $y,P_{13}$ & $2,6,8,13$ & $7,5,5,8$\\
7 & $1,2,4$ & $y,P_{13}$ & $1,4,6,8,13$ & $7,5,5,5,8$\\
8 & $1,2,3,4$ & $P_1,P_{13}$ & $1,2,4,6,8,13$ & $6,6,5,5,5,8$\\
9 & $1,2,3,4,5$ & $P_1,P_{13}$ & $1,2,4,6,8,10,13$ & $6,6,5,5,5,5,7$\\
10 & $1,2,3,4,5,6$ & $P_1,P_{13}$ & $1,2,4,6,8,10,12,13$ & $6,6,5,5,5,5,6,6$\\
\bottomrule
\end{tabular}
\end{center}
These graphs are connected and subcubic, with 27 vertices and $s=b-4$
rungs, so Lemma~\ref{lem:rungs} gives $\vp\le b$. Formula~\eqref{eq:contact}
gives the displayed distances. Successive selected $Q$-indices have absolute
distance-value difference strictly smaller than their index difference.
By the triangle inequality the same is true for any pair, so the selected
$Q$-vertices are incomparable. The selected $P_1,P_{13}$ have distance six,
and $y$ has distance seven. Cross-ray differences are at most one, except
for $y$ against a distance-five vertex and $P_1$ or $P_{13}$ against $Q_{13}$
of distance eight. The former pairs have level separation at least four.
In the latter cases $d(P_{13},Q_{13})=8$, by the closest rung $e_4$, while
$P_1,Q_{13}$ have level separation twelve. All remaining pairs with distance
difference one are nonadjacent. Thus all cross-ray comparisons are strict,
and each displayed set is a position set of size $b$.

Both rays are still $y$-geodesics, so $p_y\le2$. The graphs contain cycles,
so the first-row characterization gives $\vpm\ge2$. They realize exactly
$(2,b)$ for $6\le b\le10$. The next proposition supplies $b=2,3,4,5$;
Proposition~\ref{prop:ten} excludes every larger value.

\subsection{Intervals and exponential row growth}
\begin{proposition}\label{prop:interval}
Every integer pair $a\ge2$, $a\le b\le3a-1$ is realizable.
\end{proposition}
\begin{proof}
For $b=a$ use $K_{a+1}$. Otherwise choose $1\le r,t\le a$ with
$r+t=b+1-a$; this is possible because $2\le b+1-a\le2a$.
Take a vertex $y$, independent sets $A,D$ of sizes $r,t$, and a clique $B$
of size $a$. Add all edges from $y$ to $A$, from $A$ to $B$, and from $B$
to $D$, with no others. The graph has $b+2$ vertices. Its position numbers are
\[
\begin{array}{ll}
 p_y=a,&p_u=\max\{a+1,r+t-1,t+1\}\quad(u\in A),\\
 p_u=\max\{a,r+t-1\}\quad(u\in D),&p_u=\max\{r+a-1+t,a+t\}=b\quad(u\in B).
\end{array}
\]
For completeness, from $y$ the layers are $A,B,D$, and every earlier-layer
vertex is comparable with every later-layer vertex. From $u\in A$ the layers
are $\{y\}\cup B$ and $(A\setminus\{u\})\cup D$. A selected vertex of $B$
forbids the second layer, whereas $y$ forbids $A\setminus\{u\}$ but allows
$D$. From $u\in D$ the layers are $B$, $A\cup(D\setminus\{u\})$, and $\{y\}$;
a vertex of $B$ forbids later layers, and $y$ forbids $A$. From $u\in B$
the only nonneighbor is $y$, comparable precisely with $A$. These observations
give both the upper bounds and attaining choices in the displayed formulas.
All values lie between $a$ and $b$, with both attained.
\end{proof}

The known pairs $(a,6a-4)$ for $a\ge3$ are due to
\cite[Lemma 2.3]{position}. The exponential lower bound in
\eqref{eq:Ba} has an exact lower parameter, as follows. Set $h=a-1$, take
$a$ rays $P_0,\ldots,P_h$ attached directly to a common root $y$, and use
only the dyadic edges \eqref{eq:dyadicedges}. Put
\[
 A=4(2^h-1),\quad t_0=h+2,\quad s=A+t_0,\quad F=A+h,
\]
and give all rays length $L>s+A$. Define $C_i$ as before, and let
$x=P_{0,s}$. The potential equal to $F+1$ at $y$ and
$\min\{F+1,F_i+\dist(t,C_i)\}$ on ray $i$ is $1$-Lipschitz:
at its first vertex the untruncated value is $s+i-1=F+1+i$.
The paths used in \eqref{eq:centerdist} show that the $2^h$ centers on $P_h$
all have $x$-distance $F$.

Attach one further leaf $z$ at $y$. The graph has $a+1$ leaves, so every
observer can select at least $a$ other leaves. On the other hand the $a$
paths from $z$ through $y$ along the rays are geodesics and cover the graph;
their lengths are certified by the original level function shifted by one.
Thus $p_z=a$ and $\vpm=a$. The old target sphere is preserved. This proves
$B(a)\ge2^{a-1}$, while \eqref{eq:universal} proves the upper bound and the
existence of the maximum, since each row is nonempty and bounded in the
integers. Therefore $B(a)=2^{\Theta(a)}$.

\section{Local monophonic position and vertex visibility}
An induced path is also called a monophonic path. For $|V(G)|\ge2$, let
$m_x$ be the largest size of $S\subseteq V(G)\setminus\{x\}$ such that no
induced $x,v$-path contains another vertex of $S$, for any $v\in S$.
This localizes the universal induced-path condition of monophonic position
\cite{monophonic}. For the fan $K_1\vee P_n$ with universal vertex $w$,
$m_w=n$: every induced path from $w$ to a path vertex is its edge.
At a path vertex at most one selected path vertex lies on either side;
the only other candidate is $w$. Its neighborhood is an admissible set,
so path endpoints have value two and internal path vertices have value
three. Hence $\min_xm_x=2$, $\max_xm_x=n$ for $n\ge2$.

The existential shortest-path version already has standard terminology.
Following Roy--Di Stefano--Klav\v{z}ar--Lakshmanan
\cite{visibility}, an $x$-visibility set is a set
$S\subseteq V(G)\setminus\{x\}$ for which every $v\in S$ has some shortest
$x,v$-path meeting $S$ just at $v$. Its largest size is $v_x(G)$, and
$\vv(G)=\max_xv_x(G)$ is the vertex visibility number. We write
$\vvm(G)=\min_xv_x(G)$ for the minimum over observers. This condition does
not impose mutual visibility between two selected vertices.

\begin{proposition}\label{prop:visibility}
There is a connected cubic graph $D_t$ of order $4t+2$, for every $t\ge1$,
such that $\vv(D_t)/\vvm(D_t)\ge t/7$. Thus vertex-visibility ratios have
linear growth along a cubic family.
\end{proposition}
\begin{proof}
Start with two length-$t$ paths $P,Q$ sharing only their root $y$ and add
the rungs $P_iQ_i$ for $1\le i\le t$. Call the graph $L_t$.
Its root geodesics are the unique ray paths, so $v_y(L_t)=2$.
From $x=P_1$, the set $\{Q_1,\ldots,Q_t\}$ is visible: the paths
$P_1P_2\cdots P_iQ_i$ have length $i$, are shortest, and contain no other
selected vertex. The lower bound on length follows either by using a rung
(at least $i-1$ level-changing steps and one rung) or by going through $y$
(at least $i+1$ steps).

Take two copies of $L_t$ and join corresponding vertices $y,P_t,Q_t$.
These are exactly the degree-two vertices, so the result $D_t$ is cubic
of order $4t+2$ and is connected. Projection makes either copy isometric,
so the same target set is visible from $P_1^0$. From $y^0$, the geodesics
to $P_i^0,Q_i^0$ are unique. For $i<t$, those to $P_i^1,Q_i^1$ are also
unique: the root edge followed by the ray has length $i+1$, whereas entering
through an end edge costs at least $2t-i+1>i+1$. No rung can occur on these
root geodesics because it does not increase level.

Partition all nonobserver vertices into the four ray chains
$\{P_i^0\}_{1\le i\le t}$, $\{Q_i^0\}_{1\le i\le t}$,
$\{P_i^1\}_{1\le i<t}$, $\{Q_i^1\}_{1\le i<t}$, and the three singletons
$y^1,P_t^1,Q_t^1$. A visibility set meets each unique-geodesic chain in at
most one vertex. Hence $v_{y^0}(D_t)\le7$ and the claimed ratio follows.
The general upper bound $\vv/\vvm\le |V(G)|-1$ gives the matching order
along this family.
\end{proof}

\section{Exact checks and the remaining pair problem}
The proofs above establish the all-parameter statements. Exact finite checks
provide separate verification of the constructions. For each observer, the
geodesic order gives
\[
 p_x(G)=|V(G)|-1-\nu_x,
\]
where $\nu_x$ is the maximum matching size in the bipartite graph of all
strict comparable pairs. Indeed a chain partition gives such a matching by
linking successive vertices, and matching edges concatenate into disjoint
chains because strict order has no directed cycles; Dilworth's theorem
identifies their minimum number with the maximum antichain size.
All strict comparable pairs are necessary in this reduction, not just the
edges of $G$. Polynomial-time computability was established in
\cite[Theorem 5.7 and Corollary 5.8]{position}.

Matching and an independent recursive maximum-antichain search agree on
every observer in the five 27-vertex witnesses, all 1,834 two-ray graphs with
each ray of length one through four and every subset of allowed rungs, and
all 27,476 connected labeled graphs of orders one through six (164,031
observer values in the last collection). The two-ray upper bounds also hold
in 400 further sampled graphs with ray lengths five through fifteen.
The nine-vertex graph with edges
\[
 \{05,06,07,15,17,26,28,37,38,48,57,68\}
\]
has the independently recomputed vector
$(4,2,3,6,3,3,5,4,4)$ and hence pair $(2,6)$.

For the large constructions, breadth-first searches verify every potential
edge inequality, every specified center distance, every quotient distance
at a selected target, and every rooted covering path for $h=8,9,10$, as well
as a padded $h=8$ example. Antichain verification uses the shortest-path
directed acyclic graph: propagate a flag forward from selected vertices
along distance-increasing edges and reject any later selected vertex reached
by a flag. This checks all pairs implicitly in linear time, including targets
of different radii. The cubic graphs have orders $101454$, $204158$, $417454$
and $101488$ in these four checks. The exponential-row constructions are
checked for $2\le a\le10$, all interval examples for $2\le a\le15$, and the
visibility construction for $1\le t\le12$ and $t=20,50,100$. These finite
ranges support the proofs and make no claim of minimum order for witnesses.

The higher pair rows remain open: determine $B(a)$ exactly for $a\ge3$,
and decide whether every integer between $a$ and $B(a)$ is realizable.
Neither the exponential bounds nor the interval construction settles these
questions. Sharp leading constants in the cubic and exact-degree growth
estimates are also not determined here.

\begin{thebibliography}{9}
\bibitem{position}
M.~G.~S. Thankachy, E.~J. Thomas, U. Chandran S.~V., J. Tuite,
G. Di Stefano and G. Erskine,
\emph{On the Vertex Position Number of Graphs},
\href{https://arxiv.org/abs/2209.00359v2}{arXiv:2209.00359v2} (2022).
Numbered statements cited here refer to this version.

\bibitem{cover}
M. Dumas, F. Foucaud, A. Perez and I. Todinca,
\emph{On graphs coverable by $k$ shortest paths},
\href{https://arxiv.org/abs/2206.15088v2}{arXiv:2206.15088v2} (2023).

\bibitem{visibility}
D. Roy, G. Di Stefano, S. Klav\v{z}ar and A. Lakshmanan S,
\emph{The vertex visibility number of graphs},
\href{https://arxiv.org/abs/2510.19452v1}{arXiv:2510.19452v1} (2025).

\bibitem{dilworth}
R.~P. Dilworth, \emph{A decomposition theorem for partially ordered sets},
Ann. of Math. (2) \textbf{51} (1950), 161--166.

\bibitem{monophonic}
E.~J. Thomas, U. Chandran S.~V., J. Tuite and G. Di Stefano,
\emph{On monophonic position sets in graphs},
\href{https://arxiv.org/abs/2012.10330v4}{arXiv:2012.10330v4}.
\end{thebibliography}
\end{document}
